%% ma: L10-B13-0009
%% lop: 10
%% chuong: 5
%% bai: 13
%% dang: 10.13.3
%% loai: TL
%% muc: TH
%% dap_an: "-"
%% nguon: "chuyen-de-cac-so-dac-trung-cua-mau-so-lieu-khong-ghep-nhom-toan-10-chuong-trinh-moi.pdf (GV. Phan Nhật Linh) (P02), Bài 02 Dạng 1 Bài tập tự luận 9"
%% kien_thuc: "Lớp 10 Bài 13 (trung vị)"
%% trang_thai: chinh-thuc
%% ghi_chu: "Tự duyệt theo yêu cầu giáo viên 22/10/2026."
\begin{ex}
Số liệu sau cho biết tiền lãi hàng tháng của một cửa hàng trong năm 2021 (đơn vị: triệu đồng):
\begin{center}
\begin{tabular}{|c|c|c|c|c|c|c|c|c|c|c|c|c|}\hline
Tháng & $1$ & $2$ & $3$ & $4$ & $5$ & $6$ & $7$ & $8$ & $9$ & $10$ & $11$ & $12$\\\hline
Lãi & $12$ & $17$ & $15$ & $16$ & $18$ & $20$ & $14$ & $15$ & $18$ & $17$ & $20$ & $12$\\\hline
\end{tabular}
\end{center}
Tìm trung vị của mẫu số liệu.
\loigiai{Sắp xếp: $12$; $12$; $14$; $15$; $15$; $16$; $17$; $17$; $18$; $18$; $20$; $20$. Có $12$ giá trị nên $M_e=\dfrac{16+17}2=16{,}5$ (triệu đồng).}
\end{ex}
